\FloatBarrier
\section{模型评价与推广}
\subsection{模型优点}
\begin{enumerate}
\item SOLAR通过母体均衡权重协调不同采样频次药物对训练目标的贡献，使记录较少的药物也能参与模型优化，并与母体等权评价指标保持一致，为新药物条件下的溶解度估计提供统一的训练目标。
\item SOLAR-H将结构基线与温度响应分别建模，并通过内层分组验证选择温度增益，使条件信息的使用程度随数据证据调整，从而缓解固定温度分支带来的迁移误差，提升当前数据上的平均预测表现。
\item SOLAR-H使用低维分子描述符和树集成完成预测，在当前CPU环境下平均单次外层拟合约需零点七秒，能够以较低的计算成本完成多组参数比较，并支持交互应用中的连续条件探索与批量预测。
\item SOLAR的预测形式将结构贡献和温度贡献明确分开，可结合特征归因与响应曲线分析药物、溶剂及条件变化的作用，使预测结果便于追踪，也方便使用者在演示界面中比较不同模型的输出。
\end{enumerate}
\subsection{模型局限}
\begin{enumerate}
\item 当前数据中的独立药物母体和溶剂组合覆盖仍较有限，后续可增加不同化学类别及重复实验记录，使模型学习到更丰富的结构与条件关系，并通过新增来源的数据进一步考察预测表现的稳定性。
\item 目前输入主要由固定分子描述符与温度组成，后续可在数据条件允许时补充三维构象、酸碱条件和固态信息，以提高模型对具体实验体系的区分能力，并进一步丰富结构分支的表达方式。
\item 现有演示应用围绕数据集中已有的药物和溶剂提供选择与比较，后续可增加自定义分子输入、批量任务管理和结果检索等功能，使模型更便于融入不同规模的溶剂筛选与实验分析流程。
\end{enumerate}
\subsection{模型推广}
SOLAR-H适合用作候选药物与溶剂组合的初步排序工具，其母体均衡目标对应药物之间的平均预测质量。进一步推广可沿两条路径展开：利用来源明确的温度序列检验不同药物类别所需的收缩强度；将图表示或分子预训练嵌入结构分支\cite{molclr2022,molmerger2024}，在相同母体划分下比较其对低溶解度尾部误差的改善。前者检验条件响应的可迁移范围，后者检验更丰富表示能否补充当前描述符的信息。

\subsection{交互应用}
基于冻结模型构建的Solubility Studio采用英语为主、中文同步显示的界面，支持12种模型选择、药物与溶剂结构展示、单点预测、温度响应及多溶剂比较。两两对比使用完全相同的输入，展示预测差值、响应曲线、统一分组指标和母体误差散点；结果可导出为CSV。方法页面呈现两条算法路线和完整比较表。应用直接读取已训练模型及各自的预处理参数，将论文中的模型转化为可交互比较的预测工具。
